% Last Updated: March 1, 2024
\documentclass[a4paper, twoside]{article}
	
% Margins
\usepackage{geometry}
	\geometry{
		% showframe,
		left   = 20mm,
		right  = 20mm,
		top    = 20mm,
		bottom = 20mm
	}

% Headers and footers
\usepackage{fancyhdr}
	\pagestyle{fancy}
	\fancyhf{}
	\fancyhead[LE, RO]{\thepage}
	\fancyhead[CE]{\textsc{Cong Wen}}
	\fancyhead[CO]{\textsc\rightmark}

	\setlength{\headheight}{15pt}
	\renewcommand\headrulewidth{0pt}

% universal preamble
\input{/Users/wenc/Documents/Math/universal-pre}




\begin{document}
\title{\tbf{Title}}
\author{Cong Wen, xx/20xx}
\date{}
\maketitle
\thispagestyle{empty}
% ----------------------------------------------------------------------------------------------------

\begin{abstract}
	This is my self-use notes\footnote{Last updated on January 01, 2000} on arithmetic of number fields.

\end{abstract}
\tableofcontents


\begin{qst}
	Motivating questions
	\begin{itm}
		\item What's the inverse of $j$-invariant? Related to hypergeometric functions, Picard-Fuchs equations, Gauss-Manin connections and mirror symmetry?
		\item The Euler characteristic of arithmetic groups and Gauss-Bonnet??? See \cite[Page~255]{Brown1994} and \cite{Harder1971}.
	\end{itm}
\end{qst}

\part{Preliminaries}

\sct{Global Theory}

	\subfile{autrepn}

	\subfile{galrepn}

	\subfile{shimvar}

	\subfile{sphcvar}

    \subfile{IUTeich}


\sct{Local Theory}

	\subfile{highRam}

	\subfile{highLoc}

	\subfile{pHodge}

	\subfile{RZspace}

	\subfile{pcurve}

	\subfile{pshtuka}

	\subfile{infplace}


\part{Applications}

\sct{Arithmetic geometry}

	\subfile{ecurve}

	\subfile{mcurve}

	\subfile{modular}

	\subfile{fcurve}

	\subfile{rpoints}

	\subfile{numfield}


\sct{Arithmetic algebra}

	\subfile{invgal}

	\subfile{quatalg}


\sct{Arithmetic statistics}

	\subfile{selfdual}

	\subfile{erankone}


\sct{Arithmetic dynamics}








% ----------------------------------------------------------------------------------------------------
\bibliographystyle{amsalpha}
\bibliography{/Users/wenc/Documents/Math/universal-ref}
\end{document}

